The auxiliary gap can be introduced \emph{before} eliminating the
normalized strong root. This improves the degree-293 inner norm in
Section~\ref{sec:fourteen} to degree 219 at the same 75-operation
cost. It still uses fourteen positive unknowns. Keep normalized \(i\),
delete \(f\), replace \(y\) with \(e\), and define
\begin{align}
 Q_*&=(\Delta ic^2)^2,& d_*&=1+\Delta i^2c^4,\notag\\
 b_*&=c+Rd_*,&v_*&=cT,\notag\\
 S_*&=Q_*-1,&r_*&=b_*+S_*e,\notag\\
 p_*&=d_*v_*^2,&B_*^{\mathrm{gap}}&=2v_*r_*,\notag\\
 A_*^{\mathrm{gap}}
   &=p_*+r_*^2-Q_*S_*e^2-1,\notag\\
 K_*^{\mathrm{gap}}
   &=(A_*^{\mathrm{gap}})^2-d_*(B_*^{\mathrm{gap}})^2,
   &P_4&=N_0N_1N_kN_t.
 \label{eq:gap-root-data}
\end{align}
Here the four factors use the old transport coordinate \(z_0\).
The proposed equation is
\begin{equation}\boxed{
 \Phi_{14}^{\mathrm{gap}}
    =P_4(1+K_*^{\mathrm{gap}})-1=0.}
 \label{eq:gap14}
\end{equation}
Its evaluator computes the cross-term as \(4p_*r_*^2\), so
\(B_*^{\mathrm{gap}}\) is a proof abbreviation and needs no separate
register. Removing \(f\) and the old \(y\)-square leaves 49 prefix
operations; the new tail costs 26. The ledger is 43 multiplications
and 32 additions or subtractions.

The degree improvement comes from the exact cancellation
\begin{equation}
 A_*^{\mathrm{gap}}
  =p_*+b_*^2+(Q_*-1)e(2b_*-e)-1.
 \label{eq:gap-root-cancellation}
\end{equation}
If \(f^2=d_*\) and \(V=v_*f-b_*\), then
\begin{equation}
 A_*^{\mathrm{gap}}-B_*^{\mathrm{gap}}f
  =Q_*\bigl[V^2-(V+e)^2\bigr]+(V+e)^2-1.
 \label{eq:gap-root-remainder}
\end{equation}
Thus the new norm eliminates precisely the gap auxiliary equation.

\subsection{Restoration and the signed-domain denominators}

At a zero of \eqref{eq:gap14}, both \(P_4\) and
\(1+K_*^{\mathrm{gap}}\) are integer units. Because
\(B_*^{\mathrm{gap}}\) is even, the latter is congruent to
\(1+(A_*^{\mathrm{gap}})^2\) modulo four and cannot be \(-1\).
Hence \(P_4=1\) and \(K_*^{\mathrm{gap}}=0\). The outer unit
arguments already give \(N_0=N_1=1\) and
\(0<R<a<\Delta\). On the positive domain,
\(v_*,b_*,r_*>0\), so the denominator in
\(f=A_*^{\mathrm{gap}}/B_*^{\mathrm{gap}}\) is nonzero.
This rational root satisfies \(f^2=d_*\), hence is an integer.
Identity~\eqref{eq:gap-root-remainder}, with \(y=V+e\), restores
the normalized auxiliary equation. Replacing the multiplier by
\(i_{\mathrm{ordinary}}=\Delta i\) gives

\[
 (i_{\mathrm{ordinary}}c^2)^2
   =\Delta(f^2-1)=Q_*.
\]
The signed ordinary theorem therefore proves soundness without imposing
a sign on the recovered \(f\) or \(y\).

For the expanded domain, keep only
\(J,F,\alpha,z_0,i,s,w,\eta,\zeta>0\), and allow
\(h,T,\tau,e,\gamma\in\Z\). We must now exclude two possible
zero factors in \(B_*^{\mathrm{gap}}=2v_*r_*\).
The stronger version of the earlier outer bound is
\begin{align*}
 Q_*-b_*&=(\Delta-R)d_*-\Delta-c\\
 &\ge1+\Delta(i^2c^4-1)-c>1.
\end{align*}
Indeed \(c>2\), \(i\ge1\), and \(\Delta>0\) are already known.
Thus \(2<b_*<Q_*-1=S_*\), so

\[
 r_*=b_*+S_*e\ne0
 \qquad\text{for every integer }e.
\]

If \(T=0\), then \(v_*=p_*=B_*^{\mathrm{gap}}=0\), and
\(K_*^{\mathrm{gap}}=0\) forces \(A_*^{\mathrm{gap}}=0\).
Equivalently,
\begin{equation}
 Q_*b_*^2-(Q_*-1)(e-b_*)^2=1.
 \label{eq:gap-zero-T}
\end{equation}
Writing \(y_0=e-b_*\), the inequality \(b_*>1\) would imply
\(|y_0|\ge b_*+1\), and hence
\[
 b_*^2-1=(Q_*-1)(y_0^2-b_*^2)
       \ge(Q_*-1)(2b_*+1).
\]
This contradicts \(1<b_*<Q_*\). Therefore \(T\ne0\), both
denominator factors are nonzero, and the signed reconstruction is valid.

For completeness, the canonical positive blank witness satisfies
\(f^2=d_*\), \(P_4=1\), and \(e=y-V>0\), as proved for the
gap auxiliary coordinate. Equation~\eqref{eq:gap-root-remainder}
then gives \(A_*^{\mathrm{gap}}=B_*^{\mathrm{gap}}f\), so deleting
\(f\) yields a positive zero of \eqref{eq:gap14}. This proves
both directions of the same blank-halting and
\(\Con(\mathrm{ZFC})\) equivalence.

Nine independent positivity substitutions \eqref{eq:pos}, with the
other five coordinates left unrestricted, give a polynomial in
\(9\cdot3+5=32\) integer unknowns and 138 arithmetic operations.

\subsection{Exact degrees and the outer-norm comparison}

For the degree calculation, the following table also records weights
two on the nine positive coordinates and one on the five signed ones.
\begin{center}
\begin{tabular}{lrrrrrrrrr}
\toprule
Quantity & \(Q_*\)&\(d_*\)&\(b_*\)&\(v_*\)&\(p_*\)&\(r_*\)
 &\(A_*^{\mathrm{gap}}\)&\(B_*^{\mathrm{gap}}\)&\(P_4\)\\
\midrule
Degree &46&34&38&6&46&47&85&53&49\\
Weight &92&68&76&11&90&93&169&104&96\\
\bottomrule
\end{tabular}
\end{center}
The unique highest term of \eqref{eq:gap-root-cancellation} is
\(2Q_*eb_*\). Its square has degree 170 and weight 338, while
\(d_*(B_*^{\mathrm{gap}})^2\) has degree 140 and weight 276.
Consequently the new equation has exact degree \(49+170=219\),
and its integer conversion has exact degree \(96+338=434\).

Using \(Q_0=\beta J\), \(k_0=\eta+\zeta\), and
\(T_2=w(Q_0-F-\alpha)-z_0Q_0\), the leading form is
\begin{equation}
 (\Phi_{14}^{\mathrm{gap}})_{219}
   =32h\gamma e^2i^8k_0^{19}w^{17}s^{36}
       Q_0^{131}(Q_0-F)^2T_2.
 \label{eq:gap-root-leader}
\end{equation}
Its coefficient at
\(J^{134}h\gamma e^2i^8\eta^{19}w^{17}s^{36}z_0\)
is \(-32\beta^{134}\ne0\). This is also the weighted leading form
before the homogeneous positivity substitution, which preserves
nonvanishing.

For comparison, the outer-norm variant is
\[
 A_{\mathrm{out}}^{\mathrm{gap}}
   =P_4(A_*^{\mathrm{gap}}+1)-1,
 \qquad
 B_{\mathrm{out}}^{\mathrm{gap}}=P_4B_*^{\mathrm{gap}},
\]
\[
 \Psi_{14}^{\mathrm{gap}}
   =(A_{\mathrm{out}}^{\mathrm{gap}})^2
      -d_*(B_{\mathrm{out}}^{\mathrm{gap}})^2=0.
\]
It also costs 75 operations, but has degree 268; its 32-integer-variable
conversion has 138 operations and degree 530. The leading form is
\[
 256h^2\gamma^2e^2i^8k_0^{22}w^{22}s^{44}
       Q_0^{160}(Q_0-F)^2T_2^2.
\]
This valid form is dominated by \eqref{eq:gap14} at the same gate
count. It does not dominate the earlier 73-operation outer norm.

Its soundness proof is included to make that comparison independently
checkable. Reduction modulo \(P_4\) forces \(P_4=\pm1\), giving
the same outer bounds and \(r_*\ne0\). If \(T=0\), its equation says
that \(P_4\) times the left side of \eqref{eq:gap-zero-T} equals one.
That bracket cannot be \(-1\) modulo four because \(Q_*\) is a square;
therefore \eqref{eq:gap-zero-T} holds, an impossibility. Its root
denominator is consequently nonzero. The integer restored root now
satisfies
\[
 A_{\mathrm{out}}^{\mathrm{gap}}
    -B_{\mathrm{out}}^{\mathrm{gap}}f
  =P_4\bigl\{Q_*[V^2-(V+e)^2]+(V+e)^2\bigr\}-1=0.
\]
This restores the full normalized product, hence the signed ordinary
one. The canonical witness proves completeness as before.

Both gap root variants, their integer conversions, their live-gate
ledgers, and their algebraic checks are emitted by
\code{arithmetic/build_gap.py}. The separate proof note records the
complete coordinate substitutions and nonzero-denominator arguments.

